\section{Using Papilot}
\label{sec:using}

This section describes Papilot as an author meets it. The later sections
explain why it behaves as it does.

\subsection{Running it}

Papilot needs Node.js~22.22 or later, Git, a \TeX{} installation with
\texttt{latexmk} and \texttt{synctex}, and a signed-in agent on the host. It is
installed and started with
\begin{verbatim}
npm install -g papilot
papilot
\end{verbatim}
and then opened at \texttt{http://127.0.0.1:7274}. \emph{Add paper} registers a
\LaTeX{} Git repository on the host; later papers are added and switched from
the paper menu at the top of the sidebar. A repository must have at least one
commit and compile with \texttt{latexmk}; a custom root document is named in a
\texttt{.latex-review.json} file. To try a sample first, \texttt{papilot demo}
copies this paper into a separate repository and registers it as
\emph{Papilot paper}, without touching any existing manuscript.

\subsection{A task from start to finish}

\begin{enumerate}
  \item Select \emph{New task}, describe the revision and press Enter. The first
    line of the prompt becomes the task's title, and the task runs on the paper
    selected in the sidebar. Below the prompt, choose the agent, its model and
    its reasoning effort; Papilot remembers all three for the next task.
  \item Follow the conversation and answer the agent's questions or permission
    requests. When both execution slots are busy, further tasks queue
    (\cref{sec:agents}).
  \item When the agent finishes, Papilot builds the PDF review automatically.
    Accept, reject or comment on each change in the typeset review.
  \item Send the saved comments, or write further instructions in the
    conversation. Either starts the next round, which carries the review
    feedback with it (\cref{sec:model}).
\end{enumerate}

While the agent works, the composer shows a stop button that pauses the task.
The desktop workspace offers three views, \emph{Review}, \emph{Split} and
\emph{Chat}, and remembers the choice; on a phone the header switches between
the review and the chat. The review toolbar shows the round's progress, accepts
or rejects every undecided change at once, opens the full PDF of the round,
builds the review again, and opens it in a separate tab. \emph{Accept all} and
\emph{Reject all} act only on changes not yet decided, and one undo reverses the
whole batch.

Conversation drafts are saved per task in the browser; Enter sends, Shift+Enter
adds a line, and an input method's composition is never sent by Enter. Closing
the browser does not stop a run. After a restart of the host process, turns that
were in flight are marked interrupted and can be continued, and queued messages
are kept. If the outcome of a request is uncertain, the task offers
\emph{Check status}, and further decisions wait until it is resolved.

\subsection{Reading the paper}

\emph{PDF}, under the paper in the sidebar, shows the manuscript as it is on
disk, including accepted but uncommitted changes (\cref{sec:build}). It behaves
like a local PDF reader: continuous pages with crisp text at any zoom,
selectable text, working links and cross-references, and search across the
whole document.
\begin{itemize}
  \item \emph{Contents} lists the PDF's bookmarks, which \texttt{hyperref}
    writes for every section, highlights the section being read, and jumps to
    any entry. The same list opens from the toolbar.
  \item A page number or the arrows change page; the zoom buttons, the zoom
    menu, the keyboard and a pinch change the zoom, which is remembered.
  \item Find (Cmd+F) searches every page, with Enter and Shift+Enter stepping
    through the matches.
\end{itemize}

\subsection{Notes on the PDF}

To ask for a change directly on the PDF, select text and choose \emph{Add note},
or mark a region around a figure or an equation, then describe the change. The
selection follows the printed words exactly (\cref{sec:geometry}), each note is
settled on the source lines that produced it (\cref{sec:notes}), and every
pending note is sent to the agent as one task.

\subsection{The Git panel}

Beside the PDF, a small source-control panel shows the branch and how it
compares with its upstream, the last commit, and the uncommitted changes with
their line counts. Selecting a file shows its difference. Ticking files,
writing a message and pressing Enter commits exactly those files as they are on
disk; fetch, a fast-forward pull and push keep the branch in step with its
remote (\cref{sec:git}).

\subsection{Agents}

Each task runs on one agent: Codex through its app server, Claude through the
Claude Agent SDK, Antigravity through its command line, or Kimi Code through
ACP. For Kimi, install the CLI and run \texttt{kimi login}. Every agent uses its
own login on the host, and an agent that is not installed or signed in is listed
as unavailable. When an agent reports usage, the new-task page shows each usage
window with how much is used and when it resets (\cref{sec:usage}).

\subsection{Working from another device}

Papilot binds only to \texttt{127.0.0.1}. To reach it from a phone or another
computer, it is published on a private Tailscale network, and the server checks
the host, the browser's origin and the Tailscale identity of every remote
request (\cref{sec:security}):
\begin{verbatim}
export PAPILOT_PUBLIC_ORIGIN=https://your-mac.your-tailnet.ts.net:8443
export PAPILOT_TAILSCALE_LOGIN=your-tailscale-login
papilot service install
tailscale serve --bg --https=8443 http://127.0.0.1:7274
\end{verbatim}
A background service starts Papilot after login, restarts it if it exits and
keeps the computer from sleeping while it runs. Without Tailscale, an SSH
tunnel serves the same purpose.

\subsection{Configuration}

\begin{center}
\small
\begin{tabular}{@{}lp{0.62\linewidth}@{}}
\toprule
Variable & Purpose \\
\midrule
\texttt{PAPILOT\_PORT} & local HTTP port, $7274$ by default \\
\texttt{PAPILOT\_DATA} & state and generated example, \texttt{\textasciitilde/.papilot} by default \\
\texttt{PAPILOT\_CODEX\_BIN} & Codex executable \\
\texttt{ANTHROPIC\_API\_KEY} & API key for Claude tasks \\
\texttt{PAPILOT\_KIMI\_BIN} & Kimi Code executable \\
\texttt{PAPILOT\_AGY\_BIN} & Antigravity command line \\
\texttt{PAPILOT\_AGY\_PROXY} & HTTP proxy used only for Antigravity \\
\texttt{PAPILOT\_REVIEW\_ENGINE} & review engine directory \\
\texttt{PAPILOT\_PUBLIC\_ORIGIN} & exact browser origin for remote access \\
\texttt{PAPILOT\_TAILSCALE\_LOGIN} & Tailscale login allowed to connect \\
\bottomrule
\end{tabular}
\end{center}
